\subsection{细胞类型选择性CRE设计的活性图景建模}

公开的MPRA序列性能表包含798,064条序列记录。经过固定质量过滤、序列验证、长度过滤和精确序列去重后，全数据加载器保留了749,054条记录：627,661条训练、58,811条验证和62,582条测试。19、21和X号染色体组成验证集；7号和13号染色体组成测试集；其余符合条件的染色体组成训练集。亲本选择使用标识符和序列哈希，从不使用活性标签。

正向/反向互补平均后的Malinois预测重现了已发表的测试性能（三种细胞的平均Pearson $r=0.884$、平均Spearman $\rho=0.825$）。在相同的62,582条测试记录上，三成员残基-扩张评审集成达到平均Pearson $r=0.849$，而三成员多核评估器达到$r=0.793$（图~\ref{fig:validation}A）。它们不同的卷积结构使得一个集成可以在搜索期间评审候选，另一个则用于评估冻结的候选集。

预测不确定性由学习到的异方差方差和种子间变异共同构成，乘法缩放仅在验证数据上估计。跨细胞类型，评审模型的95\%测试区间覆盖率为0.930--0.940，跨模型评估器为0.923--0.935（图~\ref{fig:validation}B）。评审不确定性与绝对误差的Spearman相关性为$\rho=0.399$--0.448；跨模型评估器不确定性给出$\rho=0.400$--0.467（图~\ref{fig:validation}C）。拒绝曲线显示，随着高不确定性记录被逐步移除，平均绝对误差下降（图~\ref{fig:validation}D），支持不确定性引导的候选分类。

\begin{figure}[t]
  \centering
  \includegraphics[width=\textwidth]{fig2_validation.pdf}
  \caption{\textbf{预测性能和不确定性诊断。} \textbf{A}，主模型、评审模型和跨模型评估器按细胞类型划分的测试集Pearson相关系数。\textbf{B}，仅验证集方差缩放后，80\%、90\%和95\%名义水平下的观测测试覆盖率；虚线标记理想校准。\textbf{C}，预测不确定性与绝对误差之间的Spearman关联。\textbf{D}，逐步拒绝最不确定记录后的平均测试误差。所有神经集成都包含三个独立种子的成员。}
  \label{fig:validation}
\end{figure}

\subsection{SafeEdit-CRE在600个留出亲本上提升跨模型细胞类型选择性}

我们在600个留出的200-nt天然CRE亲本上评估了四种编辑方法，这些亲本是在去除所有历史亲本标识符和序列哈希后按SHA-256排序选出的。每种方法针对三个靶细胞（K562、HepG2、SK-N-SH）和四个替换预算（1、5、10、20 nt）各运行一次，产生28,800个方法级比较。随机编辑使用三次独立重复，其指标取平均，不做最优选择。候选生成后，序列被传递给未参与搜索的留出跨模型评估器。

在全部7,200个亲本--靶标--预算条件下，SafeEdit-CRE的平均跨模型特异性裕度增益为0.877，而迭代贪婪优化为0.822（配对差值0.055；亲本簇自助法95\% CI，0.040--0.071），随机匹配编辑为0.001（差值0.876；95\% CI，0.856--0.896）（表~\ref{tab:g8}；图~\ref{fig:primary}A）。跨模型裕度增益为正的条件比例，SafeEdit-CRE为95.4\%、贪婪为97.1\%，表明对贪婪的优势来自增益的幅度和分布，而非更多的正例。因此，SafeEdit-CRE的优势体现在整个条件集上更大的增益，而不是正例比例的增加。

计算量匹配的主束搜索——使用与SafeEdit-CRE相同的束宽、扩展计划和搜索预算，但仅依赖主Malinois预测器进行候选排序——的平均跨模型裕度增益为0.881。SafeEdit-CRE与这一性能高度吻合（配对差值$-0.0036$；95\% CI，$-0.0114$--0.0043；图~\ref{fig:primary}B），同时绝对靶标增益降低了0.205（95\% CI，$-0.229$--$-0.181$），跨模型集成不确定性降低了0.0209（95\% CI，$-0.0251$--$-0.0166$）（图~\ref{fig:primary}C--D）。相对于贪婪优化，SafeEdit-CRE将特异性裕度提升了0.055（95\% CI，0.040--0.071），靶标增益降低了0.141（95\% CI，$-0.177$--$-0.106$），不确定性降低了0.0153（95\% CI，$-0.0214$--$-0.0091$）。

\begin{table}[t]
\centering
\caption{600个留出亲本上的跨模型评估结果。}
\label{tab:g8}
\resizebox{\textwidth}{!}{%
\begin{tabular}{lrrrr}
\toprule
方法 & 平均裕度增益 & 平均靶标增益 & 正裕度比例(\%) & 平均不确定性 \\
\midrule
随机匹配 & 0.001 & 0.012 & 48.6 & 0.185 \\
贪婪Malinois & 0.822 & 1.803 & 97.1 & 0.329 \\
主束(计算量匹配) & 0.881 & 1.866 & 94.8 & 0.335 \\
SafeEdit-CRE & 0.877 & 1.661 & 95.4 & 0.314 \\
\bottomrule
\end{tabular}
}
\end{table}

计算量匹配的束搜索表明，仅靠搜索广度就可以恢复同等水平的预测特异性。SafeEdit-CRE改变的是从这个庞大候选池中选择哪一个：它选出的变异具有更低的集成不确定性和更小的靶标得分偏移，同时保留了特异性裕度。

\begin{figure}[t]
  \centering
  \includegraphics[width=\textwidth]{fig3_primary_endpoint.pdf}
  \caption{\textbf{600个留出天然CRE亲本上的跨模型评估。} \textbf{A}，按编辑方法划分的平均留出跨模型特异性裕度增益。\textbf{B}，SafeEdit-CRE相对于贪婪和计算量匹配主束的裕度增益配对差值。\textbf{C}，对应的绝对靶标增益配对差值。\textbf{D}，对应的跨模型集成不确定性配对差值。点表示均值；区间为100,000次重复的亲本簇自助法95\%置信区间。随机重复取平均值，从未按最优结果选择。}
  \label{fig:primary}
\end{figure}

序列域预过滤器在600个亲本中有12个（2.0\%）在至少一个条件下返回了零编辑结果。在所有汇总中保留这些结果，既保持了预设的分母，也量化了序列安全措施的实际代价。

\subsection{90个非重叠验证亲本上的设计准则基准}

为了评估SafeEdit-CRE设计不仅在跨模型评估下得分高，而且满足广泛的序列质量和模型一致性标准，我们对90条亲本序列应用了预设的九准则设计基准，这些亲本与任何开发或主要评估亲本都没有重叠（在发现6个名义亲本与早期探索性基准存在标识符级重叠后，将其排除）。
对于每个方法--亲本--靶标--预算组合，我们记录：非负主靶标增益、正主裕度增益、正评审裕度传递、链差异在域内、评审不确定性在域内、天然性在域内、GC变化在域内、同源多聚体安全性、精确编辑预算。阈值来源于验证集随机编辑，并在评估前锁定。

SafeEdit-CRE在1,080个设计中有654个（60.6\%）通过了全部九项准则，相比之下贪婪设计为419/1,080（38.8\%）、随机设计为142/1,080（13.1\%）（表~\ref{tab:g4}；图~\ref{fig:hardaudit}A）。SafeEdit--贪婪的配对差值为21.8个百分点（亲本簇自助法95\% CI，18.3--25.3个百分点）。SafeEdit--随机差值为47.4个百分点（95\% CI，43.0--51.8个百分点）。相对于贪婪的优势在一碱基预算时较小，在5--20预算时增大（图~\ref{fig:hardaudit}B--D）。

\begin{table}[t]
\centering
\caption{90个非重叠验证亲本上的九准则设计基准。}
\label{tab:g4}
\begin{tabular}{lrrrr}
\toprule
方法 & 设计数 & 通过数 & 通过率(\%) & 亲本自助法95\% CI(\%) \\
\midrule
随机匹配 & 1,080 & 142 & 13.1 & 10.6--15.8 \\
贪婪Malinois & 1,080 & 419 & 38.8 & 34.6--42.9 \\
SafeEdit-CRE & 1,080 & 654 & 60.6 & 55.6--65.3 \\
\bottomrule
\end{tabular}
\end{table}

\begin{figure}[t]
  \centering
  \includegraphics[width=\textwidth]{fig3_hard_audit.pdf}
  \caption{\textbf{九准则设计基准。} \textbf{A}，锁定的非重叠90亲本验证基准中的汇总九准则通过率。误差棒为亲本簇自助法95\%置信区间。\textbf{B--D}，分别为K562、HepG2和SK-N-SH按编辑预算分层的通过率。}
  \label{fig:hardaudit}
\end{figure}

21.8个百分点的差值衡量了在候选生成期间强制执行评审、不确定性和序列域准则的收益。对贪婪设计进行事后筛选，能够通过相同预设准则的变异体明显更少。

\subsection{细胞类型和编辑预算效应}

跨模型选择的收益因靶细胞类型和编辑预算而异（图~\ref{fig:budgettradeoff}）。对于K562，SafeEdit--贪婪的跨模型裕度差值从1个编辑的$-0.016$进展到20个编辑的0.242；对于HepG2，从$-0.004$进展到0.294。在SK-N-SH中，各预算下的差值范围为$-0.078$到$-0.033$，表明活性图景中存在细胞类型特异性的异质性。

与计算量匹配的主束相比，SafeEdit-CRE的差值在所有分层中都保持在紧密的等效范围内：K562为0.004--0.031；HepG2为$-0.030$--0.004；SK-N-SH为$-0.033$--0.010，证实了汇总的裕度等效性在靶细胞和编辑预算间均成立。

\begin{figure}[t]
  \centering
  \includegraphics[width=\textwidth]{fig5_budget_tradeoff.pdf}
  \caption{\textbf{细胞类型和编辑预算效应。} \textbf{A--B}，贪婪、主束和SafeEdit-CRE搜索在三个靶细胞上汇总的跨模型特异性裕度增益和不确定性随编辑预算的变化。\textbf{C--D}，SafeEdit-CRE相对于贪婪或计算量匹配主束的按靶标--预算的平均特异性裕度增益差值。A--B面板省略了随机编辑，因其汇总裕度增益接近零。线条和热图单元格为亲本上的描述性均值。编辑预算以nt（核苷酸替换）为单位。}
  \label{fig:budgettradeoff}
\end{figure}

这一模式与一致性重排序随着合理编辑路径数增长而变得越来越有价值是一致的。在一个替换时，三种定向策略常常都被约束在亲本序列周围的同一小邻域内；在5--20个替换时，跨模型评审过滤掉了那些在主预测器下看起来有利、但在留出集成中传递较差的编辑。SK-N-SH的结果提示，设计阈值应按细胞类型校准，而不是在细胞模型间不加改变地转移。

\subsection{完整重搜索消融识别关键组件}

为了确定哪些SafeEdit-CRE组件驱动了观察到的行为，我们在180个亲本上以七种配置重新运行了完整搜索（15,120行）。每种配置生成自己的候选序列；过滤器不是从共享候选表中追溯移除的。跨模型推理在每次搜索后都为每个候选重新计算。

完整SafeEdit-CRE的平均跨模型裕度增益为0.890，靶标增益为1.705，不确定性为0.312（图~\ref{fig:ablation}）。匹配的主束的裕度增益为0.890（配对完整--主束差值0.0002；95\% CI，$-0.0145$--0.0151），靶标增益为1.897（差值$-0.192$；95\% CI，$-0.244$--$-0.143$），不确定性为0.334（差值$-0.0212$；95\% CI，$-0.0283$--$-0.0140$）。

移除天然性重排序后，跨模型裕度增益升至0.979、靶标增益升至2.017，不确定性升至0.343。相对于此配置，完整SafeEdit-CRE将不确定性降低了0.0305（95\% CI，0.0218--0.0394），同时裕度调整了0.089（95\% CI，0.075--0.102）、靶标增益调整了0.312（95\% CI，0.266--0.360）。移除序列域预过滤器消除了全部48个不可行行，并通过放宽GC和同源多聚体限制将裕度增益提升至0.910。链和不确定性惩罚的汇总效应较小，而无评审配置则精确重现了主束结果。

\begin{figure}[t]
  \centering
  \includegraphics[width=\textwidth]{fig6_ablation.pdf}
  \caption{\textbf{180个亲本上的完整重搜索消融。} \textbf{A}，相对于完整SafeEdit-CRE的跨模型特异性裕度增益变化。\textbf{B}，对应的绝对靶标增益变化。\textbf{C}，选定候选的平均跨模型集成不确定性。\textbf{D}，不可行搜索行数。每种配置均独立重运行。天然性感知重排序是不确定性降低的主要驱动因素，而预过滤器的执行则以裕度换取合成可行性。}
  \label{fig:ablation}
\end{figure}

天然性感知重排序对不确定性降低的贡献最大，而序列域预过滤器决定了合成可行性。移除评审器精确重现了主束结果，表明评审模型改变了返回的序列，而不仅仅是在选择后对其做注释。

\subsection{Tier A候选文库及代表性示例}

从90亲本验证基准的3,240个设计中，我们导出了一个分层候选文库用于实验筛选。74条Tier A序列由SafeEdit-CRE生成，通过了全部九项准则，满足靶标和裕度要求，无合成风险标记，且位于或接近主--评审帕累托前沿（补充表~S9；补充图~S3）。Tier B包含1,141条备选序列，它们通过了核心活性和跨模型检查，但被帕累托支配或轻度偏离某条序列质量准则；其余2,025条Tier C序列保留用于失效模式分析。得到的74序列Tier A面板覆盖全部三种靶细胞类型和四种编辑预算，提供了一个紧凑的首轮报告基因筛选集，而Tier B保留了足够广度用于更大规模的后续MPRA文库。

为说明SafeEdit-CRE设计的性质，补充表~S10展示了三个代表性Tier A候选——K562、HepG2和SK-N-SH各一个——均在5编辑预算下生成并达到帕累托前沿状态且无合成风险标记。每个示例报告了亲本和编辑后的200-nt序列、编辑位置、主模型和评审特异性裕度增益、主靶细胞活性增益、跨模型预测不确定性的降低、GC含量、天然性以及同源多聚体状态。三者均显示正的主靶标增益和正的评审裕度传递，同时跨模型不确定性降低。编辑位置附近的局部序列变化被记录为可检验的序列语法假说，供实验随访。
